\begin{theorem}[真实输入约束下的 40 状态乘积下界]\label{thm:real-input-40-product-lower-bound}
令 \(X\subset\{0,1\}^{\ZZ}\) 为黄金均值移位（禁止相邻子词 \(11\)），并考虑两路真实输入 \((x,y)\in X\times X\)。
在如下自然类中，\S\ref{app:real-input-40-kernel} 的 \(40\) 状态机制 \(\Sigma=Q\times\{0,1\}^2\) 达到状态数最小：
\begin{quote}
\(\mathcal{C}\)：所有一阶（\(1\)-step）确定性机制对象，其路径同时生成三元序列 \((x,y,z)\)，满足
\begin{enumerate}
  \item \((x,y)\in X\times X\)（输入合法性作为内生约束编码于状态）；
  \item \(z\) 为在线归一化输出（与同步核 \(K_{\mathrm{sync}}\) 实现同一顺序函数）；
  \item 观测层显式携带两路输入符号 \((x_k,y_k)\)（不允许将 \((x_k,y_k)\) 预先合并为更大字母表以规避合法性记忆）。
\end{enumerate}
\end{quote}
则任意 \(M\in\mathcal{C}\) 的状态数满足 \(|\mathrm{States}(M)|\ge 40\)，且本附录构造取到等号。
\end{theorem}

\begin{proof}
证明分三步给出两个不可替代的下界并取乘积。
\par
\textbf{输入合法性记忆的 4 状态下界。}
黄金均值移位 \(X\) 为 \(1\)-step SFT：合法性仅依赖上一位是否为 \(1\)。
因此任意一阶确定性机制在单路输入上必须区分两类上下文（上一位为 \(0\) 与上一位为 \(1\)）；否则在同一状态下可允许的下一符号集合将取并集，从而必引入被禁的延拓 \(11\)。
两路独立输入时，上下文对 \((p_x,p_y)\in\{0,1\}^2\) 必须在状态中可读出，故至少需要 \(2\times 2=4\) 个互异输入上下文态。
\par
\textbf{归一化残余的 10 状态下界。}
固定任一输入上下文 \((p_x,p_y)\)，在允许输入对的子语言上，机制仍需实现同一在线归一化顺序函数（即同步核 \(K_{\mathrm{sync}}\) 的 Mealy 语义）。
由定理 \ref{thm:sync-kernel-mealy-minimality}，任何实现该顺序函数的确定性 Mealy 实现其可达状态数至少为 \(10\)。
因此在每个固定输入上下文下，机制至少需要 \(10\) 个可区分的归一化残余态。
\par
\textbf{乘积下界。}
在 \(\mathcal{C}\) 的一阶且显式输入口径下，输入合法性上下文决定可行输入边集，而归一化残余决定输出与下一残余；二者都必须由当前状态同时确定，且不可相互替代。
因此状态数至少为二者组合的笛卡尔积规模 \(4\cdot 10=40\)。
而 \(\Sigma=Q\times\{0,1\}^2\) 的构造正是把这两类信息取直积并以确定性方式耦合（见注 \ref{rem:real-input-40-product-kernel}），故达到最小。证毕。
\end{proof}

\endinput
